% appendix_diagnostics.tex
% Self-contained appendix for ASME conference paper (asmeconf class).
% Usage: place % appendix_diagnostics.tex
% Self-contained appendix for ASME conference paper (asmeconf class).
% Usage: place % appendix_diagnostics.tex
% Self-contained appendix for ASME conference paper (asmeconf class).
% Usage: place % appendix_diagnostics.tex
% Self-contained appendix for ASME conference paper (asmeconf class).
% Usage: place \input{appendix_diagnostics} after \appendix in the main document.

\section{Pipeline Data Diagnostics}\label{appendix:diagnostics}

This appendix summarizes the data volumes, distributions, and graph
statistics produced by the DFM knowledge-extraction pipeline described
in the main text.  All counts reflect the final, merged extraction
output after deduplication and quality filtering.

%% ----------------------------------------------------------------
%% Pipeline Stage Summary
%% ----------------------------------------------------------------

\begin{table}[t]
\caption{Pipeline stage summary from raw scrape to knowledge graph.}
\label{tab:pipeline-summary}
\centering
\begin{tabular}{lr}
\toprule
Stage & Count \\
\midrule
Raw threads scraped           & 12{,}805 \\
Threads after GPT relevance filter & 639 \\
Posts with answers             & 1{,}399 \\
Posts sent to LLM extraction   & 1{,}936 \\
Unique threads                 & 911 \\
Unique authors                 & 1{,}138 \\
Knowledge frames extracted     & 3{,}879 \\
Avg.\ frames per post         & 3.22 \\
Embedding vectors (3 collections) & 11{,}637 \\
Graph nodes                    & 7{,}864 \\
Structural edges               & 8{,}776 \\
Semantic edges                 & 4{,}372 \\
\bottomrule
\end{tabular}
\end{table}

The pipeline begins by scraping question threads from five Stack
Exchange sites via the public API.  A GPT-based relevance filter
retains only threads pertaining to design-for-manufacturability,
reducing the corpus from 12{,}805 to 639 threads.  Each retained
thread is normalized to plain text, and both questions and answers
(1{,}936 posts total) are sent to a structured-output LLM call that
extracts typed \emph{knowledge frames}.  The resulting frames,
posts, authors, and threads form a provenance graph augmented with
semantically inferred inter-frame edges.

%% ----------------------------------------------------------------
%% Frame Type Distribution
%% ----------------------------------------------------------------

\begin{table}[t]
\caption{Distribution of extracted knowledge frames by type.}
\label{tab:frame-types}
\centering
\begin{tabular}{lrr}
\toprule
Frame Type & Count & \% \\
\midrule
heuristic   & 1{,}401 & 36.1 \\
observation & 1{,}345 & 34.7 \\
risk        & 428     & 11.0 \\
principle   & 396     & 10.2 \\
comparison  & 149     &  3.8 \\
workaround  & 145     &  3.7 \\
case        &  15     &  0.4 \\
\bottomrule
\end{tabular}
\end{table}

Each frame is classified into one of seven substantive types (an
eighth, \emph{other}, is reserved for edge cases and was not observed
in the final dataset).  \textbf{Heuristics} capture personal
practices, rules of thumb, or directed advice together with a
rationale.  \textbf{Observations} record empirical, factual, or
definitional claims that carry no explicit action or causal mechanism.
\textbf{Risks} link a set of conditions to a failure mode or
operational limitation.  \textbf{Principles} provide causal
explanations of \emph{why} a phenomenon occurs.
\textbf{Comparisons} weigh two or more alternatives against each
other.  \textbf{Workarounds} describe non-standard fixes for stated
limitations.  \textbf{Cases} are specific real-world job or incident
narratives.

%% ----------------------------------------------------------------
%% Scope Distribution
%% ----------------------------------------------------------------

\begin{table}[t]
\caption{Distribution of extracted frames by scope domain.}
\label{tab:scope-dist}
\centering
\begin{tabular}{lrr}
\toprule
Scope Domain & Count & \% \\
\midrule
design\_geometry     & 1{,}166 & 30.1 \\
machining\_process   & 581     & 15.0 \\
other                & 512     & 13.2 \\
machine\_capability  & 507     & 13.1 \\
quality              & 235     &  6.1 \\
programming          & 225     &  5.8 \\
tooling              & 225     &  5.8 \\
material             & 165     &  4.3 \\
shop\_operations     & 141     &  3.6 \\
cost                 & 122     &  3.1 \\
\bottomrule
\end{tabular}
\end{table}

Scope domains indicate the primary technical area a frame addresses.
The dominance of \emph{design\_geometry} (30.1\%) reflects the
pipeline's focus on DFM: many frames concern feature tolerances,
wall thicknesses, hole placement, and other geometry-driven
constraints.  \emph{Machining\_process} and
\emph{machine\_capability} together account for a further 28.1\%,
covering feeds, speeds, tool paths, and machine-specific limitations.
The \emph{other} category (13.2\%) captures frames whose scope spans
multiple domains or falls outside the predefined taxonomy.

%% ----------------------------------------------------------------
%% Epistemic Stance Distribution
%% ----------------------------------------------------------------

\begin{table}[t]
\caption{Epistemic stance distribution across extracted frames.}
\label{tab:epistemic-dist}
\centering
\begin{tabular}{lrr}
\toprule
Epistemic Stance & Count & \% \\
\midrule
direct\_neutral           & 2{,}619 & 67.5 \\
tentative\_subjective     & 1{,}022 & 26.4 \\
hedged\_emphatic          & 127     &  3.3 \\
absolute\_authoritative   &  94     &  2.4 \\
narrative\_implied        &  17     &  0.4 \\
\bottomrule
\end{tabular}
\end{table}

Epistemic stance captures the confidence and authority with which the
original author presents a claim.  The majority of frames (67.5\%)
are \emph{direct\_neutral}---stated matter-of-factly without hedging
or strong emphasis.  \emph{Tentative\_subjective} frames (26.4\%)
include qualifiers such as ``I think,'' ``usually,'' or ``in my
experience.''  \emph{Hedged\_emphatic} frames combine hedging
language with strong conviction, while
\emph{absolute\_authoritative} frames present claims as definitive
rules.  \emph{Narrative\_implied} captures knowledge embedded in
storytelling rather than explicit assertion.

%% ----------------------------------------------------------------
%% Applicability Distribution
%% ----------------------------------------------------------------

\begin{table}[t]
\caption{Applicability levels assigned to extracted frames.}
\label{tab:applicability-dist}
\centering
\begin{tabular}{lrr}
\toprule
Applicability & Count & \% \\
\midrule
context\_dependent & 2{,}724 & 70.2 \\
universal          & 1{,}115 & 28.7 \\
situational        &    40  &  1.0 \\
\bottomrule
\end{tabular}
\end{table}

Applicability indicates how broadly a frame generalizes.
\emph{Context\_dependent} frames (70.2\%) apply only under specific
conditions---such as a particular material, machine type, or feature
geometry---and require an accompanying \texttt{example\_context}
field.  \emph{Universal} frames (28.7\%) hold across a wide range of
conditions without significant qualification.
\emph{Situational} frames (1.0\%) are tied to a highly specific,
one-off scenario.

%% ----------------------------------------------------------------
%% Edge Type Distribution
%% ----------------------------------------------------------------

\begin{table}[t]
\caption{Edge types in the provenance knowledge graph.}
\label{tab:edge-types}
\centering
\begin{tabular}{llr}
\toprule
Edge Type & Category & Count \\
\midrule
POST\_CONTAINS\_FRAME   & Structural & 3{,}879 \\
THREAD\_CONTAINS\_POST  & Structural & 1{,}936 \\
POST\_AUTHORED\_BY      & Structural & 1{,}936 \\
SE\_ANSWERS             & Structural & 1{,}025 \\
\addlinespace
SPECIFIES               & Semantic   & 2{,}159 \\
MITIGATES               & Semantic   & 1{,}218 \\
CAUSAL\_LINK            & Semantic   &   637 \\
CONTRADICTS             & Semantic   &   358 \\
\bottomrule
\end{tabular}
\end{table}

The knowledge graph contains two categories of edges.
\textbf{Structural edges} encode provenance: which post a frame was
extracted from, which thread contains the post, who authored it, and
the question--answer relationship between Stack Exchange posts.
\textbf{Semantic edges} are inferred by a separate LLM pass that
classifies the relationship between pairs of frames sharing similar
subjects or embeddings.  \emph{SPECIFIES} indicates that one frame
adds detail or constraints to another.  \emph{MITIGATES} links a
frame (typically a heuristic or workaround) to a risk it addresses.
\emph{CAUSAL\_LINK} connects frames sharing a cause--effect
relationship.  \emph{CONTRADICTS} flags frames that present
conflicting claims, enabling downstream systems to surface
disagreements.

%% ----------------------------------------------------------------
%% Source Site Distribution
%% ----------------------------------------------------------------

\begin{table}[t]
\caption{Distribution of posts sent to extraction by Stack Exchange site.}
\label{tab:source-sites}
\centering
\begin{tabular}{lrr}
\toprule
Stack Exchange Site & Posts & \% \\
\midrule
Engineering        & 1{,}345 & 69.5 \\
Electronics        &   276   & 14.3 \\
Robotics           &   174   &  9.0 \\
Woodworking        &   106   &  5.5 \\
Mechanics          &    35   &  1.8 \\
\addlinespace
\textbf{Total}     & \textbf{1{,}936} & \textbf{100} \\
\bottomrule
\end{tabular}
\end{table}

Posts were drawn from five Stack Exchange sites selected for their
coverage of manufacturing, machining, and related engineering
knowledge.  The \emph{Engineering} site contributes the majority of
DFM-relevant content (69.5\%), while \emph{Electronics} and
\emph{Robotics} provide complementary material on electromechanical
integration and automation constraints.  \emph{Woodworking}
contributes subtractive-manufacturing knowledge that transfers to
CNC contexts, and \emph{Mechanics} provides foundational physics
and materials content.
 after \appendix in the main document.

\section{Pipeline Data Diagnostics}\label{appendix:diagnostics}

This appendix summarizes the data volumes, distributions, and graph
statistics produced by the DFM knowledge-extraction pipeline described
in the main text.  All counts reflect the final, merged extraction
output after deduplication and quality filtering.

%% ----------------------------------------------------------------
%% Pipeline Stage Summary
%% ----------------------------------------------------------------

\begin{table}[t]
\caption{Pipeline stage summary from raw scrape to knowledge graph.}
\label{tab:pipeline-summary}
\centering
\begin{tabular}{lr}
\toprule
Stage & Count \\
\midrule
Raw threads scraped           & 12{,}805 \\
Threads after GPT relevance filter & 639 \\
Posts with answers             & 1{,}399 \\
Posts sent to LLM extraction   & 1{,}936 \\
Unique threads                 & 911 \\
Unique authors                 & 1{,}138 \\
Knowledge frames extracted     & 3{,}879 \\
Avg.\ frames per post         & 3.22 \\
Embedding vectors (3 collections) & 11{,}637 \\
Graph nodes                    & 7{,}864 \\
Structural edges               & 8{,}776 \\
Semantic edges                 & 4{,}372 \\
\bottomrule
\end{tabular}
\end{table}

The pipeline begins by scraping question threads from five Stack
Exchange sites via the public API.  A GPT-based relevance filter
retains only threads pertaining to design-for-manufacturability,
reducing the corpus from 12{,}805 to 639 threads.  Each retained
thread is normalized to plain text, and both questions and answers
(1{,}936 posts total) are sent to a structured-output LLM call that
extracts typed \emph{knowledge frames}.  The resulting frames,
posts, authors, and threads form a provenance graph augmented with
semantically inferred inter-frame edges.

%% ----------------------------------------------------------------
%% Frame Type Distribution
%% ----------------------------------------------------------------

\begin{table}[t]
\caption{Distribution of extracted knowledge frames by type.}
\label{tab:frame-types}
\centering
\begin{tabular}{lrr}
\toprule
Frame Type & Count & \% \\
\midrule
heuristic   & 1{,}401 & 36.1 \\
observation & 1{,}345 & 34.7 \\
risk        & 428     & 11.0 \\
principle   & 396     & 10.2 \\
comparison  & 149     &  3.8 \\
workaround  & 145     &  3.7 \\
case        &  15     &  0.4 \\
\bottomrule
\end{tabular}
\end{table}

Each frame is classified into one of seven substantive types (an
eighth, \emph{other}, is reserved for edge cases and was not observed
in the final dataset).  \textbf{Heuristics} capture personal
practices, rules of thumb, or directed advice together with a
rationale.  \textbf{Observations} record empirical, factual, or
definitional claims that carry no explicit action or causal mechanism.
\textbf{Risks} link a set of conditions to a failure mode or
operational limitation.  \textbf{Principles} provide causal
explanations of \emph{why} a phenomenon occurs.
\textbf{Comparisons} weigh two or more alternatives against each
other.  \textbf{Workarounds} describe non-standard fixes for stated
limitations.  \textbf{Cases} are specific real-world job or incident
narratives.

%% ----------------------------------------------------------------
%% Scope Distribution
%% ----------------------------------------------------------------

\begin{table}[t]
\caption{Distribution of extracted frames by scope domain.}
\label{tab:scope-dist}
\centering
\begin{tabular}{lrr}
\toprule
Scope Domain & Count & \% \\
\midrule
design\_geometry     & 1{,}166 & 30.1 \\
machining\_process   & 581     & 15.0 \\
other                & 512     & 13.2 \\
machine\_capability  & 507     & 13.1 \\
quality              & 235     &  6.1 \\
programming          & 225     &  5.8 \\
tooling              & 225     &  5.8 \\
material             & 165     &  4.3 \\
shop\_operations     & 141     &  3.6 \\
cost                 & 122     &  3.1 \\
\bottomrule
\end{tabular}
\end{table}

Scope domains indicate the primary technical area a frame addresses.
The dominance of \emph{design\_geometry} (30.1\%) reflects the
pipeline's focus on DFM: many frames concern feature tolerances,
wall thicknesses, hole placement, and other geometry-driven
constraints.  \emph{Machining\_process} and
\emph{machine\_capability} together account for a further 28.1\%,
covering feeds, speeds, tool paths, and machine-specific limitations.
The \emph{other} category (13.2\%) captures frames whose scope spans
multiple domains or falls outside the predefined taxonomy.

%% ----------------------------------------------------------------
%% Epistemic Stance Distribution
%% ----------------------------------------------------------------

\begin{table}[t]
\caption{Epistemic stance distribution across extracted frames.}
\label{tab:epistemic-dist}
\centering
\begin{tabular}{lrr}
\toprule
Epistemic Stance & Count & \% \\
\midrule
direct\_neutral           & 2{,}619 & 67.5 \\
tentative\_subjective     & 1{,}022 & 26.4 \\
hedged\_emphatic          & 127     &  3.3 \\
absolute\_authoritative   &  94     &  2.4 \\
narrative\_implied        &  17     &  0.4 \\
\bottomrule
\end{tabular}
\end{table}

Epistemic stance captures the confidence and authority with which the
original author presents a claim.  The majority of frames (67.5\%)
are \emph{direct\_neutral}---stated matter-of-factly without hedging
or strong emphasis.  \emph{Tentative\_subjective} frames (26.4\%)
include qualifiers such as ``I think,'' ``usually,'' or ``in my
experience.''  \emph{Hedged\_emphatic} frames combine hedging
language with strong conviction, while
\emph{absolute\_authoritative} frames present claims as definitive
rules.  \emph{Narrative\_implied} captures knowledge embedded in
storytelling rather than explicit assertion.

%% ----------------------------------------------------------------
%% Applicability Distribution
%% ----------------------------------------------------------------

\begin{table}[t]
\caption{Applicability levels assigned to extracted frames.}
\label{tab:applicability-dist}
\centering
\begin{tabular}{lrr}
\toprule
Applicability & Count & \% \\
\midrule
context\_dependent & 2{,}724 & 70.2 \\
universal          & 1{,}115 & 28.7 \\
situational        &    40  &  1.0 \\
\bottomrule
\end{tabular}
\end{table}

Applicability indicates how broadly a frame generalizes.
\emph{Context\_dependent} frames (70.2\%) apply only under specific
conditions---such as a particular material, machine type, or feature
geometry---and require an accompanying \texttt{example\_context}
field.  \emph{Universal} frames (28.7\%) hold across a wide range of
conditions without significant qualification.
\emph{Situational} frames (1.0\%) are tied to a highly specific,
one-off scenario.

%% ----------------------------------------------------------------
%% Edge Type Distribution
%% ----------------------------------------------------------------

\begin{table}[t]
\caption{Edge types in the provenance knowledge graph.}
\label{tab:edge-types}
\centering
\begin{tabular}{llr}
\toprule
Edge Type & Category & Count \\
\midrule
POST\_CONTAINS\_FRAME   & Structural & 3{,}879 \\
THREAD\_CONTAINS\_POST  & Structural & 1{,}936 \\
POST\_AUTHORED\_BY      & Structural & 1{,}936 \\
SE\_ANSWERS             & Structural & 1{,}025 \\
\addlinespace
SPECIFIES               & Semantic   & 2{,}159 \\
MITIGATES               & Semantic   & 1{,}218 \\
CAUSAL\_LINK            & Semantic   &   637 \\
CONTRADICTS             & Semantic   &   358 \\
\bottomrule
\end{tabular}
\end{table}

The knowledge graph contains two categories of edges.
\textbf{Structural edges} encode provenance: which post a frame was
extracted from, which thread contains the post, who authored it, and
the question--answer relationship between Stack Exchange posts.
\textbf{Semantic edges} are inferred by a separate LLM pass that
classifies the relationship between pairs of frames sharing similar
subjects or embeddings.  \emph{SPECIFIES} indicates that one frame
adds detail or constraints to another.  \emph{MITIGATES} links a
frame (typically a heuristic or workaround) to a risk it addresses.
\emph{CAUSAL\_LINK} connects frames sharing a cause--effect
relationship.  \emph{CONTRADICTS} flags frames that present
conflicting claims, enabling downstream systems to surface
disagreements.

%% ----------------------------------------------------------------
%% Source Site Distribution
%% ----------------------------------------------------------------

\begin{table}[t]
\caption{Distribution of posts sent to extraction by Stack Exchange site.}
\label{tab:source-sites}
\centering
\begin{tabular}{lrr}
\toprule
Stack Exchange Site & Posts & \% \\
\midrule
Engineering        & 1{,}345 & 69.5 \\
Electronics        &   276   & 14.3 \\
Robotics           &   174   &  9.0 \\
Woodworking        &   106   &  5.5 \\
Mechanics          &    35   &  1.8 \\
\addlinespace
\textbf{Total}     & \textbf{1{,}936} & \textbf{100} \\
\bottomrule
\end{tabular}
\end{table}

Posts were drawn from five Stack Exchange sites selected for their
coverage of manufacturing, machining, and related engineering
knowledge.  The \emph{Engineering} site contributes the majority of
DFM-relevant content (69.5\%), while \emph{Electronics} and
\emph{Robotics} provide complementary material on electromechanical
integration and automation constraints.  \emph{Woodworking}
contributes subtractive-manufacturing knowledge that transfers to
CNC contexts, and \emph{Mechanics} provides foundational physics
and materials content.
 after \appendix in the main document.

\section{Pipeline Data Diagnostics}\label{appendix:diagnostics}

This appendix summarizes the data volumes, distributions, and graph
statistics produced by the DFM knowledge-extraction pipeline described
in the main text.  All counts reflect the final, merged extraction
output after deduplication and quality filtering.

%% ----------------------------------------------------------------
%% Pipeline Stage Summary
%% ----------------------------------------------------------------

\begin{table}[t]
\caption{Pipeline stage summary from raw scrape to knowledge graph.}
\label{tab:pipeline-summary}
\centering
\begin{tabular}{lr}
\toprule
Stage & Count \\
\midrule
Raw threads scraped           & 12{,}805 \\
Threads after GPT relevance filter & 639 \\
Posts with answers             & 1{,}399 \\
Posts sent to LLM extraction   & 1{,}936 \\
Unique threads                 & 911 \\
Unique authors                 & 1{,}138 \\
Knowledge frames extracted     & 3{,}879 \\
Avg.\ frames per post         & 3.22 \\
Embedding vectors (3 collections) & 11{,}637 \\
Graph nodes                    & 7{,}864 \\
Structural edges               & 8{,}776 \\
Semantic edges                 & 4{,}372 \\
\bottomrule
\end{tabular}
\end{table}

The pipeline begins by scraping question threads from five Stack
Exchange sites via the public API.  A GPT-based relevance filter
retains only threads pertaining to design-for-manufacturability,
reducing the corpus from 12{,}805 to 639 threads.  Each retained
thread is normalized to plain text, and both questions and answers
(1{,}936 posts total) are sent to a structured-output LLM call that
extracts typed \emph{knowledge frames}.  The resulting frames,
posts, authors, and threads form a provenance graph augmented with
semantically inferred inter-frame edges.

%% ----------------------------------------------------------------
%% Frame Type Distribution
%% ----------------------------------------------------------------

\begin{table}[t]
\caption{Distribution of extracted knowledge frames by type.}
\label{tab:frame-types}
\centering
\begin{tabular}{lrr}
\toprule
Frame Type & Count & \% \\
\midrule
heuristic   & 1{,}401 & 36.1 \\
observation & 1{,}345 & 34.7 \\
risk        & 428     & 11.0 \\
principle   & 396     & 10.2 \\
comparison  & 149     &  3.8 \\
workaround  & 145     &  3.7 \\
case        &  15     &  0.4 \\
\bottomrule
\end{tabular}
\end{table}

Each frame is classified into one of seven substantive types (an
eighth, \emph{other}, is reserved for edge cases and was not observed
in the final dataset).  \textbf{Heuristics} capture personal
practices, rules of thumb, or directed advice together with a
rationale.  \textbf{Observations} record empirical, factual, or
definitional claims that carry no explicit action or causal mechanism.
\textbf{Risks} link a set of conditions to a failure mode or
operational limitation.  \textbf{Principles} provide causal
explanations of \emph{why} a phenomenon occurs.
\textbf{Comparisons} weigh two or more alternatives against each
other.  \textbf{Workarounds} describe non-standard fixes for stated
limitations.  \textbf{Cases} are specific real-world job or incident
narratives.

%% ----------------------------------------------------------------
%% Scope Distribution
%% ----------------------------------------------------------------

\begin{table}[t]
\caption{Distribution of extracted frames by scope domain.}
\label{tab:scope-dist}
\centering
\begin{tabular}{lrr}
\toprule
Scope Domain & Count & \% \\
\midrule
design\_geometry     & 1{,}166 & 30.1 \\
machining\_process   & 581     & 15.0 \\
other                & 512     & 13.2 \\
machine\_capability  & 507     & 13.1 \\
quality              & 235     &  6.1 \\
programming          & 225     &  5.8 \\
tooling              & 225     &  5.8 \\
material             & 165     &  4.3 \\
shop\_operations     & 141     &  3.6 \\
cost                 & 122     &  3.1 \\
\bottomrule
\end{tabular}
\end{table}

Scope domains indicate the primary technical area a frame addresses.
The dominance of \emph{design\_geometry} (30.1\%) reflects the
pipeline's focus on DFM: many frames concern feature tolerances,
wall thicknesses, hole placement, and other geometry-driven
constraints.  \emph{Machining\_process} and
\emph{machine\_capability} together account for a further 28.1\%,
covering feeds, speeds, tool paths, and machine-specific limitations.
The \emph{other} category (13.2\%) captures frames whose scope spans
multiple domains or falls outside the predefined taxonomy.

%% ----------------------------------------------------------------
%% Epistemic Stance Distribution
%% ----------------------------------------------------------------

\begin{table}[t]
\caption{Epistemic stance distribution across extracted frames.}
\label{tab:epistemic-dist}
\centering
\begin{tabular}{lrr}
\toprule
Epistemic Stance & Count & \% \\
\midrule
direct\_neutral           & 2{,}619 & 67.5 \\
tentative\_subjective     & 1{,}022 & 26.4 \\
hedged\_emphatic          & 127     &  3.3 \\
absolute\_authoritative   &  94     &  2.4 \\
narrative\_implied        &  17     &  0.4 \\
\bottomrule
\end{tabular}
\end{table}

Epistemic stance captures the confidence and authority with which the
original author presents a claim.  The majority of frames (67.5\%)
are \emph{direct\_neutral}---stated matter-of-factly without hedging
or strong emphasis.  \emph{Tentative\_subjective} frames (26.4\%)
include qualifiers such as ``I think,'' ``usually,'' or ``in my
experience.''  \emph{Hedged\_emphatic} frames combine hedging
language with strong conviction, while
\emph{absolute\_authoritative} frames present claims as definitive
rules.  \emph{Narrative\_implied} captures knowledge embedded in
storytelling rather than explicit assertion.

%% ----------------------------------------------------------------
%% Applicability Distribution
%% ----------------------------------------------------------------

\begin{table}[t]
\caption{Applicability levels assigned to extracted frames.}
\label{tab:applicability-dist}
\centering
\begin{tabular}{lrr}
\toprule
Applicability & Count & \% \\
\midrule
context\_dependent & 2{,}724 & 70.2 \\
universal          & 1{,}115 & 28.7 \\
situational        &    40  &  1.0 \\
\bottomrule
\end{tabular}
\end{table}

Applicability indicates how broadly a frame generalizes.
\emph{Context\_dependent} frames (70.2\%) apply only under specific
conditions---such as a particular material, machine type, or feature
geometry---and require an accompanying \texttt{example\_context}
field.  \emph{Universal} frames (28.7\%) hold across a wide range of
conditions without significant qualification.
\emph{Situational} frames (1.0\%) are tied to a highly specific,
one-off scenario.

%% ----------------------------------------------------------------
%% Edge Type Distribution
%% ----------------------------------------------------------------

\begin{table}[t]
\caption{Edge types in the provenance knowledge graph.}
\label{tab:edge-types}
\centering
\begin{tabular}{llr}
\toprule
Edge Type & Category & Count \\
\midrule
POST\_CONTAINS\_FRAME   & Structural & 3{,}879 \\
THREAD\_CONTAINS\_POST  & Structural & 1{,}936 \\
POST\_AUTHORED\_BY      & Structural & 1{,}936 \\
SE\_ANSWERS             & Structural & 1{,}025 \\
\addlinespace
SPECIFIES               & Semantic   & 2{,}159 \\
MITIGATES               & Semantic   & 1{,}218 \\
CAUSAL\_LINK            & Semantic   &   637 \\
CONTRADICTS             & Semantic   &   358 \\
\bottomrule
\end{tabular}
\end{table}

The knowledge graph contains two categories of edges.
\textbf{Structural edges} encode provenance: which post a frame was
extracted from, which thread contains the post, who authored it, and
the question--answer relationship between Stack Exchange posts.
\textbf{Semantic edges} are inferred by a separate LLM pass that
classifies the relationship between pairs of frames sharing similar
subjects or embeddings.  \emph{SPECIFIES} indicates that one frame
adds detail or constraints to another.  \emph{MITIGATES} links a
frame (typically a heuristic or workaround) to a risk it addresses.
\emph{CAUSAL\_LINK} connects frames sharing a cause--effect
relationship.  \emph{CONTRADICTS} flags frames that present
conflicting claims, enabling downstream systems to surface
disagreements.

%% ----------------------------------------------------------------
%% Source Site Distribution
%% ----------------------------------------------------------------

\begin{table}[t]
\caption{Distribution of posts sent to extraction by Stack Exchange site.}
\label{tab:source-sites}
\centering
\begin{tabular}{lrr}
\toprule
Stack Exchange Site & Posts & \% \\
\midrule
Engineering        & 1{,}345 & 69.5 \\
Electronics        &   276   & 14.3 \\
Robotics           &   174   &  9.0 \\
Woodworking        &   106   &  5.5 \\
Mechanics          &    35   &  1.8 \\
\addlinespace
\textbf{Total}     & \textbf{1{,}936} & \textbf{100} \\
\bottomrule
\end{tabular}
\end{table}

Posts were drawn from five Stack Exchange sites selected for their
coverage of manufacturing, machining, and related engineering
knowledge.  The \emph{Engineering} site contributes the majority of
DFM-relevant content (69.5\%), while \emph{Electronics} and
\emph{Robotics} provide complementary material on electromechanical
integration and automation constraints.  \emph{Woodworking}
contributes subtractive-manufacturing knowledge that transfers to
CNC contexts, and \emph{Mechanics} provides foundational physics
and materials content.
 after \appendix in the main document.

\section{Pipeline Data Diagnostics}\label{appendix:diagnostics}

This appendix summarizes the data volumes, distributions, and graph
statistics produced by the DFM knowledge-extraction pipeline described
in the main text.  All counts reflect the final, merged extraction
output after deduplication and quality filtering.

%% ----------------------------------------------------------------
%% Pipeline Stage Summary
%% ----------------------------------------------------------------

\begin{table}[t]
\caption{Pipeline stage summary from raw scrape to knowledge graph.}
\label{tab:pipeline-summary}
\centering
\begin{tabular}{lr}
\toprule
Stage & Count \\
\midrule
Raw threads scraped           & 12{,}805 \\
Threads after GPT relevance filter & 639 \\
Posts with answers             & 1{,}399 \\
Posts sent to LLM extraction   & 1{,}936 \\
Unique threads                 & 911 \\
Unique authors                 & 1{,}138 \\
Knowledge frames extracted     & 3{,}879 \\
Avg.\ frames per post         & 3.22 \\
Embedding vectors (3 collections) & 11{,}637 \\
Graph nodes                    & 7{,}864 \\
Structural edges               & 8{,}776 \\
Semantic edges                 & 4{,}372 \\
\bottomrule
\end{tabular}
\end{table}

The pipeline begins by scraping question threads from five Stack
Exchange sites via the public API.  A GPT-based relevance filter
retains only threads pertaining to design-for-manufacturability,
reducing the corpus from 12{,}805 to 639 threads.  Each retained
thread is normalized to plain text, and both questions and answers
(1{,}936 posts total) are sent to a structured-output LLM call that
extracts typed \emph{knowledge frames}.  The resulting frames,
posts, authors, and threads form a provenance graph augmented with
semantically inferred inter-frame edges.

%% ----------------------------------------------------------------
%% Frame Type Distribution
%% ----------------------------------------------------------------

\begin{table}[t]
\caption{Distribution of extracted knowledge frames by type.}
\label{tab:frame-types}
\centering
\begin{tabular}{lrr}
\toprule
Frame Type & Count & \% \\
\midrule
heuristic   & 1{,}401 & 36.1 \\
observation & 1{,}345 & 34.7 \\
risk        & 428     & 11.0 \\
principle   & 396     & 10.2 \\
comparison  & 149     &  3.8 \\
workaround  & 145     &  3.7 \\
case        &  15     &  0.4 \\
\bottomrule
\end{tabular}
\end{table}

Each frame is classified into one of seven substantive types (an
eighth, \emph{other}, is reserved for edge cases and was not observed
in the final dataset).  \textbf{Heuristics} capture personal
practices, rules of thumb, or directed advice together with a
rationale.  \textbf{Observations} record empirical, factual, or
definitional claims that carry no explicit action or causal mechanism.
\textbf{Risks} link a set of conditions to a failure mode or
operational limitation.  \textbf{Principles} provide causal
explanations of \emph{why} a phenomenon occurs.
\textbf{Comparisons} weigh two or more alternatives against each
other.  \textbf{Workarounds} describe non-standard fixes for stated
limitations.  \textbf{Cases} are specific real-world job or incident
narratives.

%% ----------------------------------------------------------------
%% Scope Distribution
%% ----------------------------------------------------------------

\begin{table}[t]
\caption{Distribution of extracted frames by scope domain.}
\label{tab:scope-dist}
\centering
\begin{tabular}{lrr}
\toprule
Scope Domain & Count & \% \\
\midrule
design\_geometry     & 1{,}166 & 30.1 \\
machining\_process   & 581     & 15.0 \\
other                & 512     & 13.2 \\
machine\_capability  & 507     & 13.1 \\
quality              & 235     &  6.1 \\
programming          & 225     &  5.8 \\
tooling              & 225     &  5.8 \\
material             & 165     &  4.3 \\
shop\_operations     & 141     &  3.6 \\
cost                 & 122     &  3.1 \\
\bottomrule
\end{tabular}
\end{table}

Scope domains indicate the primary technical area a frame addresses.
The dominance of \emph{design\_geometry} (30.1\%) reflects the
pipeline's focus on DFM: many frames concern feature tolerances,
wall thicknesses, hole placement, and other geometry-driven
constraints.  \emph{Machining\_process} and
\emph{machine\_capability} together account for a further 28.1\%,
covering feeds, speeds, tool paths, and machine-specific limitations.
The \emph{other} category (13.2\%) captures frames whose scope spans
multiple domains or falls outside the predefined taxonomy.

%% ----------------------------------------------------------------
%% Epistemic Stance Distribution
%% ----------------------------------------------------------------

\begin{table}[t]
\caption{Epistemic stance distribution across extracted frames.}
\label{tab:epistemic-dist}
\centering
\begin{tabular}{lrr}
\toprule
Epistemic Stance & Count & \% \\
\midrule
direct\_neutral           & 2{,}619 & 67.5 \\
tentative\_subjective     & 1{,}022 & 26.4 \\
hedged\_emphatic          & 127     &  3.3 \\
absolute\_authoritative   &  94     &  2.4 \\
narrative\_implied        &  17     &  0.4 \\
\bottomrule
\end{tabular}
\end{table}

Epistemic stance captures the confidence and authority with which the
original author presents a claim.  The majority of frames (67.5\%)
are \emph{direct\_neutral}---stated matter-of-factly without hedging
or strong emphasis.  \emph{Tentative\_subjective} frames (26.4\%)
include qualifiers such as ``I think,'' ``usually,'' or ``in my
experience.''  \emph{Hedged\_emphatic} frames combine hedging
language with strong conviction, while
\emph{absolute\_authoritative} frames present claims as definitive
rules.  \emph{Narrative\_implied} captures knowledge embedded in
storytelling rather than explicit assertion.

%% ----------------------------------------------------------------
%% Applicability Distribution
%% ----------------------------------------------------------------

\begin{table}[t]
\caption{Applicability levels assigned to extracted frames.}
\label{tab:applicability-dist}
\centering
\begin{tabular}{lrr}
\toprule
Applicability & Count & \% \\
\midrule
context\_dependent & 2{,}724 & 70.2 \\
universal          & 1{,}115 & 28.7 \\
situational        &    40  &  1.0 \\
\bottomrule
\end{tabular}
\end{table}

Applicability indicates how broadly a frame generalizes.
\emph{Context\_dependent} frames (70.2\%) apply only under specific
conditions---such as a particular material, machine type, or feature
geometry---and require an accompanying \texttt{example\_context}
field.  \emph{Universal} frames (28.7\%) hold across a wide range of
conditions without significant qualification.
\emph{Situational} frames (1.0\%) are tied to a highly specific,
one-off scenario.

%% ----------------------------------------------------------------
%% Edge Type Distribution
%% ----------------------------------------------------------------

\begin{table}[t]
\caption{Edge types in the provenance knowledge graph.}
\label{tab:edge-types}
\centering
\begin{tabular}{llr}
\toprule
Edge Type & Category & Count \\
\midrule
POST\_CONTAINS\_FRAME   & Structural & 3{,}879 \\
THREAD\_CONTAINS\_POST  & Structural & 1{,}936 \\
POST\_AUTHORED\_BY      & Structural & 1{,}936 \\
SE\_ANSWERS             & Structural & 1{,}025 \\
\addlinespace
SPECIFIES               & Semantic   & 2{,}159 \\
MITIGATES               & Semantic   & 1{,}218 \\
CAUSAL\_LINK            & Semantic   &   637 \\
CONTRADICTS             & Semantic   &   358 \\
\bottomrule
\end{tabular}
\end{table}

The knowledge graph contains two categories of edges.
\textbf{Structural edges} encode provenance: which post a frame was
extracted from, which thread contains the post, who authored it, and
the question--answer relationship between Stack Exchange posts.
\textbf{Semantic edges} are inferred by a separate LLM pass that
classifies the relationship between pairs of frames sharing similar
subjects or embeddings.  \emph{SPECIFIES} indicates that one frame
adds detail or constraints to another.  \emph{MITIGATES} links a
frame (typically a heuristic or workaround) to a risk it addresses.
\emph{CAUSAL\_LINK} connects frames sharing a cause--effect
relationship.  \emph{CONTRADICTS} flags frames that present
conflicting claims, enabling downstream systems to surface
disagreements.

%% ----------------------------------------------------------------
%% Source Site Distribution
%% ----------------------------------------------------------------

\begin{table}[t]
\caption{Distribution of posts sent to extraction by Stack Exchange site.}
\label{tab:source-sites}
\centering
\begin{tabular}{lrr}
\toprule
Stack Exchange Site & Posts & \% \\
\midrule
Engineering        & 1{,}345 & 69.5 \\
Electronics        &   276   & 14.3 \\
Robotics           &   174   &  9.0 \\
Woodworking        &   106   &  5.5 \\
Mechanics          &    35   &  1.8 \\
\addlinespace
\textbf{Total}     & \textbf{1{,}936} & \textbf{100} \\
\bottomrule
\end{tabular}
\end{table}

Posts were drawn from five Stack Exchange sites selected for their
coverage of manufacturing, machining, and related engineering
knowledge.  The \emph{Engineering} site contributes the majority of
DFM-relevant content (69.5\%), while \emph{Electronics} and
\emph{Robotics} provide complementary material on electromechanical
integration and automation constraints.  \emph{Woodworking}
contributes subtractive-manufacturing knowledge that transfers to
CNC contexts, and \emph{Mechanics} provides foundational physics
and materials content.
